\chapter*{Start here: from zero}
\addcontentsline{toc}{chapter}{Start here: from zero}
\markboth{Start here}{Start here}

This book starts at the point where school algebra ends, and ends at the
mathematics inside a working machine-learning library. If you can solve two
equations in two unknowns, you can begin. This chapter explains what ``linear
algebra'' is for, fixes the notation, and gives the single-variable calculus
needed for Chapter~11, so that nothing later arrives unannounced.

\section*{Why this subject exists}

A spreadsheet row of numbers is a \emph{vector}. A spreadsheet is a \emph{matrix}.
A rule that turns one row into another by multiplying and adding --- a
photo filter, a currency conversion, one layer of a neural network --- is a
\emph{linear map}. Linear algebra is the study of these objects, and it is useful
because the same few ideas (solve a system, find the best fit, find the
important directions) show up in every field that handles data.

\section*{A system of equations, from school}

Solve $2x - y = 0$ and $-x + 2y = 3$ by substitution. The first equation gives
$y=2x$. Putting that into the second: $-x+2(2x)=3$, so $3x=3$, so $x=1$ and $y=2$.
Check: $2(1)-2=0$ and $-1+2(2)=3$. Chapter~1 solves this same system two other
ways and then solves systems with hundreds of unknowns, where substitution would
be hopeless.

\section*{The notation, in one page}

\begin{center}
\begin{tabular}{lP{9.2cm}}
\toprule
\thead{Symbol} & \thead{Meaning}\\
\midrule
$\mathbb R$ & the real numbers, i.e.\ the number line\\
$\mathbb R^{n}$ & lists of $n$ real numbers; $\mathbb R^{2}$ is the plane, $\mathbb R^{3}$ is space\\
$\mathbf x$ & a vector, written as a column, such as $\begin{pmatrix}1\\2\end{pmatrix}$\\
$A$ & a matrix, a rectangular table; $A$ is $m\times n$ if it has $m$ rows and $n$ columns\\
$a_{ij}$ & the entry of $A$ in row $i$, column $j$\\
$A^{\mathsf T}$ & the transpose: swap rows and columns\\
$\sum_{i=1}^{n}a_i$ & $a_1+a_2+\dots+a_n$\\
$\in$ & ``is an element of''; $\mathbf x\in\mathbb R^{3}$ says $\mathbf x$ is a list of three numbers\\
$\forall$, $\exists$ & ``for every'', ``there exists''\\
\bottomrule
\end{tabular}
\end{center}

\section*{A calculus primer for Chapter~11}

You will not need calculus until Chapter~11, but it is better to meet it here.
The \emph{derivative} $f'(x)$ of a function $f$ is the slope of its graph at $x$:
how fast $f$ changes when $x$ is nudged. For example $f(x)=x^2$ has $f'(x)=2x$;
at $x=3$ the slope is $6$, so a nudge of $0.01$ raises $f$ by about $0.06$
(indeed $3.01^2-3^2=0.0601$). The \emph{chain rule} handles a function inside a
function: if $f(x)=(3x+1)^2$ then $f'(x)=2(3x+1)\cdot 3=18x+6$, the derivative of
the outside evaluated at the inside, times the derivative of the inside. You can
confirm by expanding $(3x+1)^2=9x^2+6x+1$ and differentiating directly.
Chapter~11 extends exactly these two ideas to functions of many variables.

\section*{How the books of this series fit together}

This is Book~3 of eight, and it is the mathematical spine of the data half of the
series.

\begin{center}
\begin{tabular}{P{2.4cm}P{9.0cm}}
\toprule
\thead{Book} & \thead{Connection}\\
\midrule
2 Discrete maths & Teaches how to read and write proofs; this book states each theorem and then proves it in that style. Chapter~12 here builds on its finite probability.\\
4 Machine learning & Cites this book by chapter: least squares (Chapter~8), PCA (Chapter~10), gradients and backpropagation (Chapters~11 and~14), losses (Chapter~15).\\
6 Quant research & Uses the statistics of Chapter~13 and the covariance ideas of Chapters~9--10.\\
\bottomrule
\end{tabular}
\end{center}

\section*{How to use the boxes in each chapter}

Every chapter opens with \emph{After this chapter you can}, listing concrete
skills. From Chapter~2, a \emph{Before you start} box names the earlier
chapters it relies on and gives a one-sentence recap, so you can tell whether to
reread something. Each chapter ends with a \emph{Hand derivation} section: do it
on paper before you trust the code.
